\section{Applications}
\label{sec:applications}

We first transfer the total-depth gaps to trivially perfect
completion, then specify the vertex-query search conventions and their
hardness transfers. The YES deletion sets and NO separator bounds also
give hardness for deletion and labeling objectives. Established
automata and hypergraph reductions transfer the parameter gaps,
while the linear treewidth gaps yield exponential lower bounds
for junction-tree table size.

Several application routes are summarized below. The separator
promise means the line-graph YES deletion and NO separator bounds
from Section~\ref{sec:sseh}; the objective definitions and the
individual transfers follow in the indicated sections.
\[
\begin{aligned}
 \text{total depth}
 &\xrightarrow{\text{Lemma~\ref{lem:tp-fill-depth}}}
   \text{trivially perfect completion},\\
 \text{line-graph separator promise}
 &\xrightarrow{\text{Section~\ref{subsec:deletion-labeling}}}
   \text{deletion and label sizes},\\
 \text{treewidth}
 &\xrightarrow{\text{Sections~\ref{subsec:established-transfers} and~\ref{subsec:total-state-space}}}
   \text{hypergraph widths and table size}.
\end{aligned}
\]

\subsection{Trivially perfect completion by number of added edges}

The total-depth transfer must account for the input's original edges.

The \emph{closure} of a rooted forest joins each vertex to all its
proper ancestors. A graph is \emph{trivially perfect} if and only if
it is such a closure~\cite[Section~2]{HogemoEtAl2021}. Let $f(J)$ be
the minimum number of edges
added to $J$ to obtain a trivially perfect graph, and recall that
$\totd(J)$ is its minimum total elimination depth.
\begin{lemma}[Completion identity {\cite[Theorem~8]{HogemoEtAl2021}}]
\label{lem:tp-fill-depth}
For every graph $J$,
\begin{equation}\label{eq:tp-fill-depth}
 f(J)=\totd(J)-|V(J)|-|E(J)|.
\end{equation}
\end{lemma}

For any $C\geq1$, a factor-$C$ completion of cost $z$ satisfies
\begin{equation}\label{eq:tp-completion-guarantee}
 f(J)\leq z\leq C f(J).
\end{equation}
Adding the known offset $|V(J)|+|E(J)|$ cannot worsen the
multiplicative guarantee: the offset is added once to the returned
cost, while a factor-$C$ comparison allows $C$ times the offset. Thus
\begin{equation}\label{eq:tp-adjusted-cost}
\begin{aligned}
 \totd(J)
  &\overset{\eqref{eq:tp-fill-depth}}{=}f(J)+|V(J)|+|E(J)|
 \overset{\eqref{eq:tp-completion-guarantee}}{\leq}z+|V(J)|+|E(J)|\\
 &\overset{\eqref{eq:tp-completion-guarantee}}{\leq}C f(J)+|V(J)|+|E(J)|
 \leq C\bigl(f(J)+|V(J)|+|E(J)|\bigr)
  \overset{\eqref{eq:tp-fill-depth}}{=}C\totd(J),
\end{aligned}
\end{equation}
where the final inequality uses $C\ge1$.
Thus Corollary~\ref{cor:total-depth-bounded-degree} also gives NP-hardness
within every fixed factor below $63/62$ and UGC hardness
below $8/7$ for trivially perfect completion on graphs of bounded
maximum degree. From a supplied trivially perfect supergraph $K$,
construct a generating forest in polynomial time by choosing a
universal vertex as the root of each connected component, deleting
it, and recursing on the remaining components. Such a vertex exists
by the closure characterization, which also holds for induced
subgraphs by restricting their generating forests. The resulting
forest has closure $K$ and total depth $|E(K)|+|V(K)|$, so the
adjusted cost also supplies the forest needed by the total-depth
approximation transfer.

\Needspace{6\baselineskip}
\begin{theorem}\label{thm:tp-completion-sse}
For every fixed $C\ge1$, approximating the minimum number of edges
added to obtain a trivially perfect graph within factor $C$ is
SSE-hard on graphs of maximum degree bounded by a constant depending
only on $C$.
\end{theorem}

\begin{proof}
Apply~\eqref{eq:tp-adjusted-cost} to the SSE hardness of total
depth in Corollary~\ref{cor:total-depth-bounded-degree}.
\end{proof}

Deciding whether at most $k$ edges suffice to complete a graph to a
trivially perfect graph is NP-hard. Chain Completion asks for the
fewest edges added across a fixed bipartition to make the
neighborhoods on one side linearly ordered by inclusion. Its
NP-hardness~\cite{Yannakakis1981} and the polynomial reduction to
trivially perfect completion~\cite[Lemma~5.4]{BliznetsEtAl2016}
give this decision hardness.
Bliznets,
Cygan, Komosa, Mach, and Pilipczuk show that, under the Exponential
Time Hypothesis (ETH), for some
fixed $c\ge1$ there is no
$2^{O(k^{1/4}/\log^c k)}n^{O(1)}$-time algorithm for this completion
problem, where $k$ is the number of added edges and $n$ is the input
order~\cite[Theorem~1.1]{BliznetsEtAl2016}.
The SSE hardness of \emph{chordal} Minimum Fill-In at every fixed factor,
whose target graphs have no induced cycle of length at least four,
is already known~\cite[Theorem~1.2]{APW}. Restricting the
target class to trivially perfect graphs changes the optimum, so
that theorem alone does not prove the statement here.
% Archived with the alternative incidence-based treedepth proof.
% The incidence construction of Wu, Austrin, Pitassi, and Liu
% ~\cite[Sections~4.3 and~5]{APW} yields a bipartite input construction.
% We derive the trivially perfect completion hardness for these inputs
% below. Let $n=|V(G)|$ and take $R_{\mathrm{inc}}=nm$, as in
% Appendix~\ref{app:incidence-proof}. Then $N'=|V(G')|=R_{\mathrm{inc}}n+m
% =m(n^2+1)$, and the output has at most
% $2R_{\mathrm{inc}}m=2nm^2$ edges. Its YES total depth is at most
% $4\varepsilon mN'$, since its treedepth is at most $4\varepsilon m$.
% Appendix~\ref{app:incidence-proof} gives an absolute constant
% $c_{\mathrm{inc}}>0$ for this construction. In the NO case,
% Lemma~\ref{lem:balanced-by-td} gives total depth at least
% $(c_{\mathrm{inc}}/2)\sqrt{\varepsilon}\,mN'$. The completion identity
% subtracts $N'+|E(G')|$, and
% \[
%  \frac{N'+|E(G')|}{mN'}
%  \leq \frac{N'+2nm^2}{mN'}
%  =\frac1m+\frac{2n}{n^2+1}
%  \leq \frac4n,
% \]
% where the last inequality uses $m\geq n/2$, since the host is
% regular and $m>0$. For each fixed $\varepsilon$, we may exclude the
% finitely many source instances below any fixed order, so assume
% $n\geq16/(c_{\mathrm{inc}}\sqrt{\varepsilon})$. The subtraction is
% then at most $(c_{\mathrm{inc}}/4)\sqrt{\varepsilon}\,mN'$. Therefore
% the NO completion cost is at least
% $(c_{\mathrm{inc}}/4)\sqrt{\varepsilon}\,mN'$, whereas the YES cost
% is at most $4\varepsilon mN'$. Their ratio is at least
% $c_{\mathrm{inc}}/(16\sqrt{\varepsilon})$, which exceeds any fixed
% constant for a sufficiently small fixed $\varepsilon$ in the allowed
% range of the source promise and incidence construction. Hence the
% completion problem is SSE-hard to approximate within every fixed
% constant on bipartite inputs as well.

\subsection{Vertex-query search}
\label{subsec:vertex-query-search}

On a connected graph $G$ with $n=|V(G)|$, a vertex query either finds the target or
identifies the component containing it after the queried vertex is
removed. Under the confirming convention, the target must itself be
queried, even if it is the only remaining candidate. With unit query
costs, the minimum sum of query counts over all targets is $\totd(G)$;
with a uniform target distribution, the minimum expected cost is
$\totd(G)/|V(G)|$.

Some search models instead stop without a query once the target is
uniquely identified. If $I(G)$ is the minimum sum of query counts
under this identification convention, then
\begin{equation}\label{eq:search-identification-offset}
 I(G)\leq\totd(G)\leq I(G)+|V(G)|.
\end{equation}
A confirming strategy can stop early, giving the first inequality.
Conversely, add a final query of the identified target on each outcome
where it has not yet been queried. This costs at most one extra query
per target and gives the second inequality. Thus $\totd(G)/|V(G)|$
and $I(G)/|V(G)|$ differ by at most one, but need not be equal:
on the three-vertex path their values are $5/3$ and $1$, respectively.

The identification approximation of~\cite{Szyfelbein2025GraphSearch}
also gives an $O(\sqrt{\log n})$-approximation for the uniform
confirming objective. Add a final confirming query where necessary.
On a connected graph with at least two vertices every strategy makes
at least one query, so the extra expected cost of at most one is
absorbed by a constant factor. The one-vertex case is solved directly.

Our NP and UGC total-depth reductions have connected outputs. The UGC
graph $S_q(G)$ is connected whenever $G$ has an edge. The edgeless
source is always a YES instance and can be mapped to $K_1$. In the NP
construction, opposite pairs and clause pairs join $A$ to $B$ and $C$,
and every auxiliary pair has an edge between $B$ and $C$
(Section~\ref{subsec:np-total-depth}). The bounded-degree replacements
retain these joining edges and use connected sparse replacement graphs
from Lemma~\ref{lem:factor-expanders}. Thus the same constant-factor
hardness holds for average-case confirming search with uniform weights
and unit query costs, also at bounded maximum degree.

It also holds under the identification convention. Indeed, fix a target
factor $r$ strictly below the relevant threshold and first choose a
source gap with a larger ratio $r'>r$. Its YES total-depth upper bound
$Y$ is $\Theta(N^2)$ in the output order $N$, and its NO lower bound
$Z$ satisfies $Z>r'Y$. Equation~\eqref{eq:search-identification-offset}
gives identification bounds $Y$ and $Z-N$, whose gap satisfies
\[
 (Z-N)-rY>(r'-r)Y-N>0
\]
for sufficiently large source instances. The source padding used in
the reductions ensures this size condition. Dividing by $N$ preserves
the ratio for average cost. A returned strategy is converted to a
forest by adding the terminal queries and recursively using its query
as the root of the current candidate component; queries outside that
component reveal no information and may be deleted. These operations
are polynomial-time for an explicitly represented strategy.

\subsection{Deletion and labeling objectives}
\label{subsec:deletion-labeling}

\paragraph{Vertex integrity.}
For $S\subseteq V(J)$, let $\mathcal C(J-S)$ be the set of connected
components of $J-S$, and define
\[
c_J(S)=\max\bigl(\{|V(C)|:C\in\mathcal C(J-S)\}\cup\{0\}\bigr),\qquad \operatorname{vi}(J)=\min_{S\subseteq V(J)}\bigl(|S|+c_J(S)\bigr).
\]
The elementary inequality
$\td(J)\leq\operatorname{vi}(J)$ follows by putting $S$ on a chain
and its components on separate chains below it. For the layered graph
$S_q(G)$, write $n=|V(G)|$. Deleting all copies of a minimum source
vertex cover leaves $q$ separate
cliques on $\alphaI(G)$ vertices. Thus the same YES construction gives
$\operatorname{vi}(S_q(G))\overset{\text{Lem.~\ref{lem:ugc-completeness}}}{\leq}
q(n-\alphaI(G))+\alphaI(G)$,
while Theorem~\ref{thm:ugc-structural} supplies the NO lower bound.
The parameter choices in Section~\ref{subsec:ugc-td-gap} therefore
also give UGC hardness for vertex integrity below $4/3$.
For SSE hardness, let $J=\mathcal L(G^*)$ be the output in the
dense construction of Section~\ref{subsec:sse-dense-gap}, put
$N=|V(J)|$, and let $D$ consist of the output vertices corresponding
to crossing host edges. Every component of $J-D$ lies within the
vertices representing one class's internal edges. Thus
\begin{equation}\label{eq:integrity-sse-yes}
\operatorname{vi}(J)\leq |D|+c_J(D)
 \overset{\eqref{eq:sse-yes-host-bounds}}{\leq}
       3\varepsilon N+(1+\varepsilon)N/q \leq\left(3+\frac{1+\varepsilon}{2}\right)
      \varepsilon N
  \leq4\varepsilon N,
\end{equation}
using $q=\lceil2/\varepsilon\rceil$ and $\varepsilon\le1$.
In the NO case, the separator
lemma with $\rho=2/3$ gives
\begin{equation}\label{eq:application-separator-bound}
 b_{2/3}(J)\overset{\text{Lem.~\ref{lem:linegraph-separators}}}{\geq}\beta N,
\end{equation}
where
$b_{2/3}(J)$ is the minimum size of a set whose deletion leaves every
component of order at most $2N/3$, and
$\beta=(c/2)\sqrt\varepsilon$, with $c>0$ from
Lemma~\ref{lem:sse-unweighted-gap}. Choose $\varepsilon$ small enough
that the source construction in
Section~\ref{subsec:sse-dense-gap} applies, $\beta\leq1/12$, and
$\varepsilon\leq1$. For any $S\subseteq V(J)$, either $S$ is a
$2/3$-balanced separator, giving
$|S|\geq b_{2/3}(J)\overset{\eqref{eq:application-separator-bound}}{\geq}\beta N$, or
$c_J(S)>2N/3\geq\beta N$. In both cases
$|S|+c_J(S)\geq\beta N$, so
\begin{equation}\label{eq:integrity-sse-no}
 \operatorname{vi}(J)\geq\beta N.
\end{equation}
The ratio of the NO lower bound to the YES upper bound grows as
$1/\sqrt\varepsilon$, so it exceeds any fixed factor for sufficiently
small $\varepsilon$. Hence vertex integrity has no constant-factor
approximation under SSEH.

The \emph{fracture number} is
$\operatorname{fr}(J)=\min_{S\subseteq V(J)}\max\{|S|,c_J(S)\}$, and
\emph{starwidth} $\operatorname{sw}(J)$ is the minimum width of
a tree decomposition whose underlying tree is a star.
\begin{lemma}[Elementary integrity comparisons]
\label{lem:integrity-comparisons}
For every nonempty graph $J$,
\[
 \operatorname{fr}(J)\leq\operatorname{vi}(J)
 \leq2\operatorname{fr}(J),\qquad
 \operatorname{sw}(J)+1\leq\operatorname{vi}(J)
 \leq2(\operatorname{sw}(J)+1).
\]
\end{lemma}

\begin{proof}
For every $S\subseteq V(J)$,
\[
 \max\{|S|,c_J(S)\}\leq |S|+c_J(S)
 \leq2\max\{|S|,c_J(S)\}.
\]
Minimizing gives the fracture-number comparisons. A star decomposition
with center bag $S$ and one leaf bag $S\cup V(C)$ for each component $C$
of $J-S$ has width at most $|S|+c_J(S)-1$. Minimizing over $S$ gives
$\operatorname{sw}(J)+1\leq\operatorname{vi}(J)$.

Conversely, let a star decomposition of $J$ have width
$k=\operatorname{sw}(J)$, and let $S$ be its center bag. By connectedness of the bags containing a vertex,
each vertex outside $S$ occurs in only one leaf bag. Edge coverage
therefore prevents edges between vertices in different leaf bags
outside $S$. Every component of $J-S$ therefore lies in one leaf bag.
The center bag and each leaf bag have at most $k+1$
vertices, so
\begin{equation}\label{eq:star-component-bounds}
 |S|\leq k+1,\qquad c_J(S)\leq k+1.
\end{equation}
Hence
$\operatorname{vi}(J)\leq |S|+c_J(S)
\overset{\eqref{eq:star-component-bounds}}{\leq}2(k+1)$.
\end{proof}

For the same SSE instances, apply the comparisons to
$\operatorname{fr}(J)$ and $\operatorname{sw}(J)+1$:
\[
\begin{aligned}
 \textup{YES:}\quad
 \max\{\operatorname{fr}(J),\operatorname{sw}(J)+1\}
 &\overset{\text{Lem.~\ref{lem:integrity-comparisons}}}{\leq}\operatorname{vi}(J)
 \overset{\eqref{eq:integrity-sse-yes}}{\leq}4\varepsilon N,\\[2pt]
 \textup{NO:}\quad
 \min\{\operatorname{fr}(J),\operatorname{sw}(J)+1\}
 &\overset{\text{Lem.~\ref{lem:integrity-comparisons}}}{\geq}\operatorname{vi}(J)/2
 \overset{\eqref{eq:integrity-sse-no}}{\geq}\beta N/2.
\end{aligned}
\]
Since $\beta/(8\varepsilon)=c/(16\sqrt\varepsilon)$ grows without
bound as $\varepsilon$ decreases, choose $\varepsilon$ for the desired factor.
For an integer $t\ge1$, let $G_t^*$ replace each edge occurrence of
$G^*$ by $t$ parallel copies. For every $v\in V(G^*)$ and
$A\subseteq V(G^*)$,
\[
\begin{gathered}
 |E(G_t^*)|=t|E(G^*)|,\qquad
 \deg_{G_t^*}(v)=t\deg_{G^*}(v),\qquad
 |V(\mathcal L(G_t^*))|=tN,\\
 \vol_{G_t^*}(A)=t\vol_{G^*}(A),\qquad
 |\cut_{G_t^*}(A)|=t|\cut_{G^*}(A)|.
\end{gathered}
\]
Thus the normalized source promises are unchanged and the output
order grows without bound. The additive $1$ is negligible compared with the
linear YES bound $4\varepsilon N$ and is absorbed by
Lemma~\ref{lem:gap-preservation}.
Gima et al. give an $O(\log\operatorname{vi}(J))$ approximation for vertex
integrity~\cite[Theorem~5.1]{GimaEtAlIntegrity}.
Earlier NP-hardness includes vertex integrity on connected line
graphs~\cite[Corollary 4.11]{GimaEtAlIntegrity} and starwidth~\cite{VanEe2017}.

\paragraph{General and average customizable labels.}
A customizable hub labeling assigns $L(v)\subseteq V(J)$ so that
\[
 V(P)\cap L(s)\cap L(t)\ne\varnothing
 \qquad\text{for every }s\text{--}t\text{ path }P,
\]
including the zero-length path at a vertex~\cite{BlumStorandt2023}.
The labeling is \emph{hierarchical} if there is a total order
$v_1,\ldots,v_{|V(J)|}$ of $V(J)$ such that
$L(v_i)\subseteq\{v_i,\ldots,v_{|V(J)|}\}$ for every $i$.
Without this requirement we call it general. Average,
maximum, and total label size mean, respectively,
$\sum_{v\in V(J)}|L(v)|/|V(J)|$, $\max_{v\in V(J)}|L(v)|$, and
$\sum_{v\in V(J)}|L(v)|$.
\begin{lemma}[Average label size and separators {\cite[Corollary~10, full version]{BlumStorandt2023}}]
\label{lem:customizable-average-separator}
For every customizable hub labeling $L$ of a nonempty graph $J$,
\begin{equation}\label{eq:customizable-average-separator}
 \frac{\sum_{v\in V(J)}|L(v)|}{|V(J)|}
 \geq
 \frac{b_{2/3}(J)}4.
\end{equation}
\end{lemma}

Choose a fixed $\varepsilon>0$ small enough that the source
construction in Section~\ref{subsec:sse-dense-gap} applies and
$(c/2)\sqrt\varepsilon\leq1/12$ and
$\sqrt\varepsilon<c/(32C)$, where $C\ge1$ is the desired factor.
For the resulting graph $J=\mathcal L(G^*)$, order the YES forest's
vertices in postorder, so every vertex precedes its ancestors. For
this forest $F$, define
\[
 L(v)=\{u\in V(J):u\text{ is an ancestor of }v\text{ in }F\},
 \qquad v\in V(J),
\]
where a vertex is its own ancestor. This is a hierarchical labeling.
On any $s$--$t$ path, its shallowest vertex $z$ is an ancestor of
every path vertex by Lemma~\ref{lem:connected-set-ancestor}.
Thus $z\in L(s)\cap L(t)$,
so these are valid customizable labels. The YES height bound and the
NO separator bound give, respectively,
\[
 \textup{YES:}\quad\max_{v\in V(J)}|L(v)|
 \overset{\eqref{eq:sse-yes-height}}{\leq}4\varepsilon |V(J)|,
 \qquad
 \textup{NO:}\quad\frac{\sum_{v\in V(J)}|L(v)|}{|V(J)|}
 \overset{\text{Lem.~\ref{lem:customizable-average-separator}}}{\geq}\frac{b_{2/3}(J)}4
 \overset{\text{Lem.~\ref{lem:linegraph-separators}}}{\geq}
 \frac c8\sqrt\varepsilon\,|V(J)|.
\]
Consequently every constant
approximation factor is SSE-hard for average or maximum general
label size and average hierarchical label size.

For hierarchical labels, the total-depth gaps give stronger conclusions
under the other two hardness assumptions as well. The following identity
also supplies a polynomial-time decoding of labels into a forest.
\begin{lemma}[Hierarchical labels and total depth]
\label{lem:hierarchical-label-total-depth}
For every nonempty graph $J$, the minimum total size of a hierarchical
customizable hub labeling is $\totd(J)$. Its minimum average size is
$\totd(J)/|V(J)|$.
\end{lemma}
\begin{proof}
Ancestor sets in an elimination forest are valid hierarchical labels,
as proved above, and their total size is the forest's total depth.
This proves the upper bound.

Conversely, let $L$ be a hierarchical labeling and fix a compatible
total order. In each connected component choose its highest-ranked
vertex as the root, delete it, and recurse on the remaining components.
The resulting forest $F$ is an elimination forest: every graph edge
either meets the chosen root or remains within one recursive component.
Suppose $a$ is an ancestor of $v$ in $F$. The component in which $a$
was chosen contains an $a$--$v$ path $P$, all of whose vertices have
rank at most that of $a$. The cover condition gives a vertex in
$L(a)\cap L(v)\cap V(P)$. Since hierarchy permits only vertices of
rank at least that of $a$ in $L(a)$, this vertex must be $a$.
Writing $\operatorname{Anc}_F(v)$ for the ancestors of $v$ including
$v$ itself, we have proved
\begin{equation}\label{eq:hierarchical-label-ancestor-containment}
\begin{gathered}
 \operatorname{Anc}_F(v)\subseteq L(v)\quad(v\in V(J)),\\
 \totd(J)\overset{\text{definition of }\totd}{\leq}\totd(F)
 =\sum_{v\in V(J)}\operatorname{depth}_F(v)
 \overset{\eqref{eq:hierarchical-label-ancestor-containment}}{\leq}
 \sum_{v\in V(J)}|L(v)|.
\end{gathered}
\end{equation}
This proves the reverse bound. The recursion is polynomial-time.
If only the labels are supplied, recover a compatible order by
topologically sorting the constraints $v<u$ for
$u\in L(v)\setminus\{v\}$; hierarchy guarantees that these constraints
are acyclic. Dividing the total-size identity by $|V(J)|$ proves the
average-size identity.
\end{proof}

For a customizable contraction hierarchy (CCH), an elimination order adds edges
between the remaining neighbors of each eliminated vertex. Orient
the resulting edges toward the later vertex in that order.
The upward search space of $v$ is the set of vertices reachable
from $v$ along these oriented edges, including $v$ itself.
\begin{lemma}[CCH search spaces and hierarchical labels {\cite[Lemmas~2 and~4 and Corollary~5, full version]{BlumStorandt2023}}]
\label{lem:cch-hierarchical-labels}
For a fixed vertex order, the CCH upward search spaces form a
hierarchical customizable hub labeling and are contained in every
hierarchical customizable hub labeling respecting that order.
Consequently, minimizing total, average, or maximum upward vertex
search-space size over CCH orders gives the corresponding minimum
hierarchical customizable-label size.
\end{lemma}

By Lemmas~\ref{lem:cch-hierarchical-labels}
and~\ref{lem:hierarchical-label-total-depth},
the minimum total upward vertex search-space size is
also $\totd(J)$, and its minimum average is $\totd(J)/|V(J)|$.
An order gives all upward search spaces by polynomial-time graph searches;
Lemma~\ref{lem:hierarchical-label-total-depth} then decodes them
into an elimination forest.

\begin{corollary}\label{cor:hierarchical-label-hardness}
Minimum total or average hierarchical customizable-label size, and
minimum total or average CCH upward vertex search-space size, are
NP-hard to approximate within every fixed factor $1\leq r<63/62$,
Unique Games-hard within every fixed factor $1\leq r<8/7$, and
SSE-hard within every fixed factor $r\geq1$. These conclusions hold
on graphs of bounded maximum degree and on regular graphs, with the
degree depending only on the excluded factor. The SSE conclusions also
hold on bounded-degree bipartite $2$-degenerate graphs of any prescribed
fixed girth lower bound.
\end{corollary}
\begin{proof}
Apply Lemmas~\ref{lem:hierarchical-label-total-depth}
and~\ref{lem:cch-hierarchical-labels} to
Corollaries~\ref{cor:total-depth-bounded-degree},
\ref{cor:regular-hardness}, and~\ref{cor:subdivision-hardness}.
Multiplication by the known graph order preserves each approximation
factor. The polynomial decodings above also preserve these factors.
\end{proof}

\subsection{Transfers through established reductions}
\label{subsec:established-transfers}

The following consequences use existing correspondences. We state
the transformations and their losses to specify which parameter
gaps they preserve.

\paragraph{Automata and restricted cycle rank.}
The \emph{star height} of a regular language is the minimum nesting
depth of Kleene stars in a regular expression for the language.
A partial deterministic finite automaton (DFA) has an initial state,
accepting states, and at most one transition from each state per
letter, but may omit transitions. Its language consists of words labeling paths
from the initial state to an accepting state. It is
\emph{bideterministic} if it has one accepting state and remains
deterministic when its transitions are reversed and its initial and
accepting states are exchanged. It is \emph{trim} if every state lies
on such a path. A binary DFA uses a two-letter alphabet. For a digraph
$D$, the \emph{cycle rank} $\operatorname{cr}(D)$ is zero if $D$ is
acyclic, is $1+\min_{v\in V(D)}\operatorname{cr}(D-v)$ if $D$ is strongly
connected and has a directed cycle, and otherwise is the maximum
cycle rank of its strongly connected components, the maximal sets
whose vertices can reach one another by directed paths
~\cite[Sections~2--3]{Gruber2008}.

\begin{theorem}[Cycle rank and star height {\cite[proof of Theorem~12, Theorem~20, and proof of Theorem~23]{Gruber2008}}]
\label{thm:cycle-rank-star-height}
For a nonempty graph $Q$, let $D_Q$ replace each edge by its two
orientations. Then $\operatorname{cr}(D_Q)=\td(Q)-1$.
The star height of a bideterministic language equals the cycle rank
of the transition digraph of its minimal trim DFA (McNaughton's theorem).
Moreover, a bideterministic partial DFA can be transformed in
polynomial time into a binary bideterministic partial DFA preserving
the language's star height and retaining the initial and sole
accepting state.
\end{theorem}

Theorem~\ref{thm:cycle-rank-star-height} transfers our three
treedepth gaps to star height for languages represented by trim binary
bideterministic partial DFAs. We first add a universal vertex $a$
to $G$ to ensure connectivity and cancel the unit difference between
treedepth and symmetric cycle rank. Define the symmetric digraph $D$ by
\[
 V(D)=V(G)\cup\{a\},\qquad
 E(D)=\{(u,v),(v,u):uv\in E(G+a)\}.
\]
Use state set $V(D)$ and alphabet $E(D)$, with transitions
$u\xrightarrow{(u,v)}v$ for $(u,v)\in E(D)$, and take $a$ as
the initial and sole accepting state. The resulting partial DFA accepts the label sequences
of directed walks from $a$ back to $a$. It is trim because the digraph
is strongly connected, and it is bideterministic because each letter
labels only one arc. Every trim bideterministic DFA is minimal:
reverse determinism and the sole accepting state imply that a word
can be accepted from at most one state. Trimness supplies an
accepting suffix from each state, so any two states are
distinguishable by such a suffix.
Putting $a$ above an optimal forest of $G$ gives one direction of
\begin{equation}\label{eq:universal-vertex-treedepth}
 \td(G+a)=\td(G)+1.
\end{equation}
For the reverse direction, in any elimination tree of $G+a$ every
vertex is comparable with $a$, because $a$ is universal. Every
root-to-leaf path therefore contains $a$: a branch leaving its ancestor
path above $a$ would contain an incomparable vertex. Restricting the
tree to $G$ removes $a$ from each such path and lowers its height
by at least one. Gruber's symmetric-digraph correspondence now gives
\[
 \operatorname{cr}(D)
 \overset{\text{Thm.~\ref{thm:cycle-rank-star-height}}}{=}
 \td(G+a)-1
 \overset{\eqref{eq:universal-vertex-treedepth}}{=}\td(G).
\]
Apply the binary encoding in Theorem~\ref{thm:cycle-rank-star-height}
to this DFA.
Delete states that are not on an initial-to-accepting path; this
polynomial trimming preserves the language and bideterminism. The
resulting trim DFA is minimal by the preceding argument and has its
initial state as its sole accepting state.
Thus every state is reachable from the initial state and can return
to it, so the transition digraph is strongly connected. Bideterminism over two
letters bounds its indegree and outdegree by two. By
Theorem~\ref{thm:cycle-rank-star-height}, the encoded language's star
height is the cycle rank of this bounded-degree digraph.
Gruber proves cycle-rank NP-hardness
even for digraphs of maximum outdegree two and maximum total degree
four~\cite[Corollary 24]{Gruber2008}. Our consequence is the transfer
of the present approximation gaps.

\paragraph{Edge ranking.}
An \emph{edge ranking} assigns positive integer labels to edges
so that every path containing two distinct edges of the same
label also contains an edge with a larger label. Its minimum
possible largest label is the edge-ranking number
$\operatorname{er}(B)$.
\begin{lemma}[Edge ranking and line-graph treedepth {\cite{BodlaenderRankings}}]
\label{lem:edge-ranking-linegraph}
For every simple graph $B$,
\begin{equation}\label{eq:edge-ranking-linegraph}
 \operatorname{er}(B)=\td(\mathcal L(B)).
\end{equation}
\end{lemma}

Let $G^*$ be the host multigraph in Section~\ref{sec:sseh}, and put
$H=\mathcal L(G^*)$. The SSE reduction proves hardness of
$\td(H)$ at every fixed factor.
Lemma~\ref{lem:edge-ranking-linegraph} transfers treedepth hardness
of a simple graph's line graph to edge ranking.
We turn $G^*$ into a simple bipartite graph before applying this
identity. Replace every loop occurrence by an edge to a fresh pendant
vertex, keeping parallel edge occurrences distinct. This preserves
the line graph: each loop is replaced by one edge incident with its
original endpoint, and the new pendant vertices have no other
incident edges. Then subdivide each edge occurrence once, using a
fresh vertex $x_e$. Call the resulting graph $B$. Each new edge joins
an old vertex to $x_e$, so it is not a loop. Different occurrences
use different subdivision vertices, so their new edges cannot be
parallel. Thus $B$ is simple. It is bipartite,
with the old and pendant vertices on one side and the subdivision
vertices on the other.

For each edge occurrence $e=uv$ before subdivision, the two new
edges are $ux_e$ and $x_ev$. Their vertices in $\mathcal L(B)$ are
adjacent. Contracting these pairs, one pair for every edge occurrence,
recovers $H$. Thus $H$ is a minor of $\mathcal L(B)$.
Conversely, replace each vertex of
an elimination forest of $H$ by a two-vertex chain, attaching its
child subtrees below the lower copy. Every edge of $\mathcal L(B)$
joins two copies of adjacent vertices of $H$, or the two copies of
one vertex. The expanded forest is therefore valid and has at most
twice the original height, proving
\begin{equation}\label{eq:linegraph-doubling}
 \td(\mathcal L(B))\leq2\td(H),\qquad
 \td(H)\overset{\text{Lem.~\ref{lem:minor-monotonicity}}}{\leq}\td(\mathcal L(B))
 \overset{\eqref{eq:edge-ranking-linegraph}}{=}\operatorname{er}(B)
 \overset{\eqref{eq:linegraph-doubling}}{\leq}2\td(H).
\end{equation}
Applying the SSE
treedepth hardness at factor $2C$ gives hardness for edge ranking on
simple bipartite graphs at every fixed factor $C\geq1$.
Lam and Yue proved that computing an optimal edge ranking is NP-hard
at ICS 1996~\cite{LamYue1996}.

\paragraph{Hypergraph widths and binary queries.}
A hypergraph $H$ has a vertex set $V(H)$ and a collection $E(H)$ of
subsets called hyperedges; its \emph{rank} is the largest hyperedge size.
A tree decomposition consists of a tree $T$ and a bag $B_t\subseteq V(H)$
for each node $t$. Every vertex must occur in at least one bag, the nodes
whose bags contain any fixed vertex must form a connected subtree of $T$,
and each hyperedge must be contained in a bag. An edge cover of a
bag $B_t$ is a set of hyperedges $\lambda_t\subseteq E(H)$ satisfying
$B_t\subseteq\bigcup_{e\in\lambda_t}e$; these edges are also called
\emph{guards}. The \emph{generalized hypertree width} is
\[
 \operatorname{ghw}(H)
 =\min_{(T,(B_t),(\lambda_t))}\max_{t\in V(T)}|\lambda_t|,
\]
where the minimum ranges over tree decompositions with such covers.
The \emph{fractional hypertree width} $\operatorname{fhw}$
instead assigns weights $\gamma_t(e)\ge0$ to each hyperedge $e$
at each bag $B_t$, requires
\begin{equation}\label{eq:fractional-cover-condition}
 \sum_{e\in E(H):\,v\in e}\gamma_t(e)\ge1
 \qquad(t\in V(T),\ v\in B_t),
\end{equation}
and minimizes $\max_{t\in V(T)}\sum_{e\in E(H)}\gamma_t(e)$.
The ordinary \emph{hypertree width} $\operatorname{hw}(H)$ minimizes
the same objective as $\operatorname{ghw}(H)$ over rooted tree
decompositions that also satisfy the \emph{special condition}
\[
 \left(\bigcup_{e\in\lambda_t}e\right)
 \cap\left(\bigcup_{s\in V(T_t)}B_s\right)\subseteq B_t
 \qquad(t\in V(T)),
\]
where $T_t$ is the subtree rooted at $t$
~\cite[Definitions~2.4--2.7 in the full version]{GottlobEtAlHypertree}.

Marx reduces treewidth to fractional hypertree width
in~\cite[proof of Theorem~5.1]{Marx2010Fractional}. We use adjacent
pairs of copies to transfer all three widths on rank-two instances.
For a nonempty graph $G$, define $G[K_2]$ on
$V(G)\times\{1,2\}$: distinct vertices $(u,i),(v,j)$ are adjacent
if and only if $u=v$ or $uv\in E(G)$. Thus each source vertex is
replaced by an adjacent pair, and a source edge joins the two pairs
completely.
The edges of $G[K_2]$ are viewed as size-two hyperedges, with the
standard width definitions of~\cite{GottlobEtAlHypertree}.
\begin{lemma}[Rank-two width identity]
\label{lem:rank-two-width}
For every nonempty graph $G$,
\[
 \operatorname{fhw}(G[K_2])=\operatorname{ghw}(G[K_2])
 =\operatorname{hw}(G[K_2])=\operatorname{tw}(G)+1.
\]
\end{lemma}
\begin{proof}
For the upper bound, take a tree decomposition $(T,(B_t))$ of $G$
of width $\operatorname{tw}(G)$ and set, for every $t\in V(T)$,
\[
 B'_t=B_t\times\{1,2\},\qquad
 \lambda_t=\{\{(v,1),(v,2)\}:v\in B_t\}.
\]
This gives each of
the three widths an upper bound of $\operatorname{tw}(G)+1$.
The guards are contained in the bag, so the special condition for
ordinary hypertree width holds as well. In particular,
\begin{equation}\label{eq:rank-two-width-upper}
 \operatorname{hw}(G[K_2])\leq\operatorname{tw}(G)+1.
\end{equation}
For the lower bound, consider a fractional decomposition of width
$w$ with decomposition tree $T$. By the width definition,
\begin{equation}\label{eq:fractional-bag-budget}
 \sum_{e\in E(G[K_2])}\gamma_t(e)\leq w\qquad(t\in V(T)).
\end{equation}
The covering condition and the size-two hyperedges give, at every bag,
\[
 |B_t|\overset{\eqref{eq:fractional-cover-condition}}{\leq}
       \sum_{e\in E(G[K_2])}\gamma_t(e)|e\cap B_t|
 \leq2\sum_{e\in E(G[K_2])}\gamma_t(e)
       \overset{\eqref{eq:fractional-bag-budget}}{\leq}2w.
\]
Define bags
\[
 C_t=\{v\in V(G):(v,1),(v,2)\in B_t\},\qquad
 |C_t|\leq |B_t|/2\leq w\quad(t\in V(T)).
\]
Let $T_{(v,i)}$ be the
connected subtree of bags containing the copy $(v,i)$. The nodes whose
bags $C_t$ contain $v$ form
$T_{(v,1)}\cap T_{(v,2)}$. This intersection is nonempty because
the two copies are adjacent, and it is connected because both sets
are connected in $T$. If $uv\in E(G)$, its four copies form a clique,
so their sets $T_{(u,1)},T_{(u,2)},T_{(v,1)},T_{(v,2)}$ intersect pairwise.
By Lemma~\ref{lem:subtree-helly}, one bag contains
all four copies. The bags $C_t$ form
a tree decomposition of $G$ with bag size at most $w$, so
\begin{equation}\label{eq:rank-two-retained-width}
 \operatorname{tw}(G)+1\leq w.
\end{equation}
Since this holds for every fractional decomposition, the definitions give
\[
\operatorname{tw}(G)+1
 \overset{\eqref{eq:rank-two-retained-width}}{\leq}\operatorname{fhw}(G[K_2])
 \leq\operatorname{ghw}(G[K_2])\leq\operatorname{hw}(G[K_2]) \overset{\eqref{eq:rank-two-width-upper}}{\leq}\operatorname{tw}(G)+1,
\]
since fractional covers relax integral covers and hypertree
decompositions also impose the special condition. Together with the
common upper bound, this proves the three equalities.
\end{proof}

By Lemma~\ref{lem:rank-two-width}, the NP $r<12/11$, UGC $r<4/3$,
and established SSE treewidth
hardness at every fixed factor transfer even to rank-two hypergraphs.
The transfer compares $\operatorname{tw}(G)+1$, the size of a largest
bag. Adding one to both promise bounds is absorbed by
Lemma~\ref{lem:gap-preservation}, since the bounds grow linearly.
In the usual query representation, vertices are variables and
each hyperedge supplies one relation atom. Their conjunction is
a conjunctive query, and size-two hyperedges give binary atoms.
Using a distinct relation symbol for each atom therefore gives
the same structural width consequences for these queries.

\subsection{Binary junction-tree total state space}
\label{subsec:total-state-space}

The width gaps in Sections~\ref{sec:tw-transfer}
and~\ref{subsec:np-treewidth} are linear in the
output order. Exponentiating bag sizes therefore yields exponential
approximation gaps for the total size of dense binary tables.

For an $n$-vertex graph $G$ with $n\ge1$, define the total state space
(TSS) by
\[
 \operatorname{TSS}(G)=\min_{(T,(B_t)_{t\in V(T)})}
       \sum_{t\in V(T)}2^{|B_t|}.
\]
The minimum ranges over tree decompositions of $G$.
Each vertex represents a binary variable. The table for bag
$B_t$ has one entry for each of its $2^{|B_t|}$ assignments.
Thus this sum is the total number of table entries in the
corresponding junction tree.
The optimization output is a tree decomposition, with its cost
represented as a binary integer; the dense tables need not be
materialized. For a decomposition with $b$ bags this integer has
$O(n+\log(b+1))$ bits, so its cost is computable in polynomial time
from the decomposition even when the table count is exponential.
Bodlaender and Fomin identify the equivalent chordal-completion
objective and prove that, unless $\mathrm{P}=\mathrm{NP}$, no
polynomial-time algorithm approximates binary total state space within
any fixed factor~\cite[Lemma 2, Section 8, and Theorem 32]{BodlaenderFomin2005}.
\begin{lemma}[Number of bags at optimal width {\cite[proof of Lemma~4]{BodlaenderFomin2005}}]
\label{lem:optimal-width-bag-count}
Every nonempty $n$-vertex graph $G$ has a tree decomposition of width
$\operatorname{tw}(G)$ with at most $n$ bags.
\end{lemma}

Apply Lemma~\ref{lem:optimal-width-bag-count}.
Each bag then contributes at most $2^{\operatorname{tw}(G)+1}$ entries.
Conversely, every tree decomposition has a bag of size at least
$\operatorname{tw}(G)+1$, so
\begin{equation}\label{eq:binary-state-width-bounds}
 2^{\operatorname{tw}(G)+1}\leq\operatorname{TSS}(G)
       \overset{\text{Lem.~\ref{lem:optimal-width-bag-count}}}{\leq}n2^{\operatorname{tw}(G)+1}.
\end{equation}
These bounds give
NP hardness within $2^{\chi n}$ for every fixed $0<\chi<1/18$,
and UGC hardness for every fixed $0<\chi<1/6$.
Indeed, let $w_Y,w_N$ be the YES and NO treewidth bounds.
For any exponent $\chi$ below the corresponding threshold, first
choose the source error small and then the replication factor large.
This gives $w_N-w_Y\ge(\chi+\delta)n$ for some fixed $\delta>0$.
By~\eqref{eq:binary-state-width-bounds}, the NO-to-YES ratio for total
state space is at least $2^{(\chi+\delta)n}/n>2^{\chi n}$ for large
$n$. A fixed number of source copies ensures this size condition.
Appendix~\ref{app:state-space-gap} gives the width differences,
parameter choices, and threshold comparison.
Bonnet's earlier linear width gap~\cite{BonnetTreewidth} also gives
a fixed positive exponential exponent by the same calculation.

\paragraph{Equivalent formulations.}
In a \emph{vertex ranking}, every path whose distinct endpoints
have the same label has an internal vertex with a larger label.
A \emph{centered coloring} gives every nonempty connected induced
subgraph a color occurring once. The minimum numbers of
labels or colors in these two definitions equal treedepth.
These, together with symbolic elimination height,
minimum clique number of a trivially perfect completion, maximum
hierarchical customizable label size, and maximum CCH vertex
search-space size are established treedepth formulations
\cite{Reidl2014,BodlaenderRankings,BlumStorandt2023}. These formulations
are covered by the same theorem.
